\chapter{Plans, evidence and runs}
\label{ch:plans}

Two questions decide whether an agent is useful on real work: does it think
before it acts, and does it know when it is actually done? \nodwyn{}
answers both in Scheme: the \emph{plan gate} and the \emph{evidence loop}.
Runs, the batch tool for scripts and benchmarks, are built on top of both.

\section{The plan gate}

Before a turn's work, \code{(plan-gate request)} may have a plan made. With
a \textbf{judge} configured (Jev: its primitives registered and a
\code{jev-key} credential stored and bound; Chapter~\ref{ch:keys}), it goes
like this:

\begin{enumerate}
  \item The judge is asked \code{plan-needed-question}: does this request
    need a plan? Below \code{plan-needed-threshold} (0.5), the turn just
    proceeds.
  \item A read-only sub-agent labelled \code{explorer} looks at the folder
    for up to \code{plan-explore-rounds} (12) rounds (\code{fs-read},
    \code{fs-list}, \code{index-grep}, \code{fs-find}), so the plan is not
    written blind.
  \item The control model drafts a plan in a narrowed context, offered no
    tools, with the explorer's findings.
  \item The judge scores the plan against the request. Below
    \code{plan-accept-score} (0.7) the model revises it, at most
    \code{plan-max-revisions} (3) times; the best plan wins.
  \item The outcome is recorded as a \code{gate.plan} event, and the window
    pins the plan into the system prompt under \code{plan-header} for the
    rest of the turn. The explorer's findings are handed to the model that
    does the work as well, so it does not read the same files again.
\end{enumerate}

Without a judge, an ordinary interactive turn has no gate. If the judge's
question fails, the model plans anyway, unscored (``degraded''). A judge
that the guard refuses (its key is not bound) ends the turn, like any
denial.

\begin{txt}
you> :eval (list plan-needed-threshold plan-accept-score plan-max-revisions
                 plan-explore-rounds)
(0.5 0.7 3 12)
\end{txt}

All of these are plain definitions you can change in \code{init.scm}, as
are the prompts: \code{plan-system-prompt}, \code{plan-draft-prompt},
\code{plan-revise-prompt}, \code{plan-explore-system},
\code{plan-score-rubric} and \code{plan-score-levels}.

\subsection{The one invariant}

Whatever the Scheme does, the core holds one line: \textbf{when a turn
must plan, nothing irreversible runs until a non-empty plan is recorded,
and a turn that records none fails.} It enforces that a plan exists, not
who wrote it. So a project with its own way of planning can record its own
plan, as trusted project code:

\begin{scm}
(define house-plan
  """
  1. Read the files the request names, and nothing else yet.
  2. Say what you found before changing anything.
  3. Change one file at a time; run the checks after each.
  """)

(define house-base-turn nodwyn-turn)
(define (nodwyn-turn request)
  (unless (plan-recorded?)
    (record-plan (list (cons 'decision "plan") (cons 'plan house-plan)
                       (cons 'reason "the house plan"))))
  (house-base-turn request))
\end{scm}

\begin{txt}
    3 gate.plan       {Origin:gate Noul:plan Plan:1. Read the files the request names ...
\end{txt}

Note that this wraps \code{nodwyn-turn}, not \code{agent-turn}: the host
calls \code{nodwyn-turn} for every message, so this runs even in a
supervised run, and since a plan is then recorded, the default gate that
runs next has nothing left to do. It \emph{wraps} rather than replaces,
because a working habit (\file{checkloop.scm}) already wraps
\code{nodwyn-turn} to remember the request for its \code{second-look}
tool; replacing it outright would quietly switch that off. (The name is specific to \nodwyn{} so that a project's own
\code{turn}, a counter say, can never shadow it by accident.)

\subsection{The model's own plan}

Without a plan from the gate (an interactive turn with no judge), the model
writes its own with the \code{plan-write} tool: numbered steps, recorded as
\code{gate.plan} with decision \code{model} and pinned under
\code{plan-header} like any other plan. It matters more than it looks: a
model's thinking is not sent back to it, and a reply that only calls tools
says nothing, so a design worked out in thinking is worked out again at
every request until it is written down. Writing the plan again revises it.

\subsection{Progress through the plan}

The plan header asks the model to mark each step as it finishes it, with
the \code{plan-step} tool: a step's number, \code{done} or
\code{skipped}, and a note saying what shows it. The strategy answers the
call (it is an agent tool) with \code{record-step}, which records
\code{plan.step} and returns how many steps are marked and which is next;
the model reads that answer, the console's plan panel ticks the step, and
\code{:plan} in the terminal client lists the steps as \code{[x]} done,
\code{[-]} skipped or \code{[ ]} not yet, with the notes. A step the plan
does not have, or a mark in a turn with no plan, is said in the answer and
recorded nowhere. A mark is the model's word, like its report: the
transcript is the evidence. A step's number indexes the plan's numbered or
bulleted lines, which is why the planner is asked for numbered steps.

\section{The evidence loop}

When nobody reviews the work as it happens, a model's ``I'm done'' is not
enough. With \code{(unattended?)} (\code{-{}-yolo}, \code{-{}-judge}, an
unattended run) or in a supervised run, the default \code{agent-turn} runs
the \textbf{evidence loop} instead of the plain one. It ends a turn on
evidence:

\begin{itemize}
  \item A reply that stops before anything changed gets
    \code{finish-no-change-message}.
  \item A reply that stops after a change, with no test run since the last
    edit, gets \code{finish-untested-message}.
  \item At most \code{finish-reminders} (2) such reminders per turn.
  \item When a person reads the answer before anything else happens
    (\code{(reviewed?)}: a chat, or a run's turn that a checkpoint follows),
    the no-change reminder is \code{finish-no-change-reviewed-message}: make
    the change unless the request asks for none yet, a proposal to review
    first say, and after it the turn may end unchanged. An unattended run,
    and a run's last turn, still need a change.
  \item After \code{unrunnable-after} (6) test commands in a row fail to run
    at all, the model gets \code{(finish-unrunnable-message n)} once: finish
    and report the change as unverified, rather than spend the budget on a
    broken environment.
\end{itemize}

A command counts as a test run when its command line contains one of
\code{test-command-words}: \code{pytest}, \code{go test},
\code{cargo test}, \code{npm test}, \code{make test}, \code{rspec} and
twenty-one more. A test run with failures still counts as a run; only
commands that never got as far as running a test count against
\code{unrunnable-after}.

The scratch directory is not the repository. A file the model writes there,
such as its check script \code{\$TMPDIR/check.sh}, is not a change, and a
command that runs it is a test run; a check fails until the work is done, so
its failing runs never count against \code{unrunnable-after}. In a git work
tree git decides what changed, so an edit made through the shell counts and
a write that changes nothing does not. Files new since the turn began that
git does not track are named once, with
\code{(finish-new-files-message files)}, before the turn ends.

\begin{tip}
If your project runs its tests some other way, teach the loop:
\begin{scm}
(define test-command-words
  (append test-command-words '("./run_tests" "bazel test" "just check")))
\end{scm}
Otherwise the evidence loop will keep telling the model it has not tested,
however many times it runs your test script.
\end{tip}

\section{The unattended system prompt}

\code{(system-prompt repo)} returns \code{(unattended-system-prompt repo)}
when nobody is there to answer what asks (\code{(unattended?)}), and
\code{(interactive-system-prompt repo)} otherwise. The interactive one
tells the model that the user follows its work and may be asked to approve
a command, so it should say in one short line what it is about to do before
calling a tool. The unattended one is about how to work instead (understand
the request, set up a check that runs, make the whole change, validate by
running it, report honestly). Both include \code{working-rules}
(Chapter~\ref{ch:habits}).

\section{Runs: the batch tool}

\begin{sh}
nodwyn daemon                       # in one terminal, or the app
nodwyn run --task task.md --repo . --max-turns 20 --budget-wall 2h
\end{sh}

A run carries a task to its end with no conversation: a sequence of turns
driven by the daemon, not by you typing. It is a batch tool, for scripts
and benchmarks (\file{docs/runs.md}); in the console, a long request is
just a message (below). Without \code{-{}-unattended} it is a
\textbf{supervised} run, which asks you between turns. Each turn:

\begin{enumerate}
  \item sends the message \code{(run-turn-prompt task complete-line turn
    max-turns)} makes: the task, and a request to do the plan's next
    unmarked step, mark it with \code{plan-step}, and report;
  \item must plan (\code{plan-required?} is true), so \code{nodwyn-turn}
    runs the plan gate first, with or without a judge;
  \item runs the evidence loop, since nobody watches each call.
\end{enumerate}

A turn is one message and the work up to a reply with no tool call. It is
what the per-turn budgets measure, where the run checkpoint comes and where
a run can stop and resume. The \emph{plan} is the run's. The first turn
drafts it. Every later turn carries it: \code{(run-plan)} reads it from the
log with the steps marked so far, and \code{plan-carry} records it again as
this turn's (\code{gate.plan} with decision \code{carried} and the
\code{marks}), so the steps keep their numbers and the window shows what is
done under the plan. Nothing is explored, drafted or scored for it. Only
when every step is marked and the run goes on does a turn plan again, and
the planner is told the previous plan and how each step went
(\code{plan-run-request}). To set a plan aside, the model marks the steps
that no longer apply \code{skipped}, saying why. A project that wants a new
plan every turn redefines \code{plan-gate}.

Between turns you get a \textbf{run checkpoint}: a yes continues and grants
nothing else, no file, command or network permission. The run ends when:

\begin{itemize}
  \item the model's final plain reply ends with the exact line
    \code{NODWYN\_RUN\_COMPLETE}, which is the model's claim, not proof: read
    the transcript and run your checks;
  \item \code{-{}-max-turns} or \code{-{}-budget-wall} runs out;
  \item you stop it (\code{nodwyn run stop}, or Stop in the console), or a
    refusal, progress stop or provider failure does.
\end{itemize}

\code{nodwyn run resume SESSION} continues a stopped run with the same
deadline and turn count, and restores no temporary approvals. A turn that
was interrupted without a recorded outcome is never assumed to have
succeeded: the driver refuses to guess.

The console does not start or resume runs (\adr{35}). Every task it starts
is a chat, and a request that takes long is a message like any other: the
turn goes on until the model answers or a hard limit stops it, and with the
default soft limits no checkpoint interrupts it. The console shows a run
started elsewhere, and its Stop ends it; a message afterwards goes on with
the task as a chat. A \file{config.toml} that still has the \code{[tasks]}
table, which once chose how the console started a task, loads, and the
table is ignored.

A run's writes in its folder go ahead as every session's do
(\code{[approval] writes}, Chapter~\ref{ch:primitives}). Its other confirms
follow its flags: the judge answers them with \code{-{}-judge}, yes does
with \code{-{}-yolo}, and an unattended run declines them. A supervised run
with none of these flags has you to ask, so it is answered as a chat is:
\code{[approval] commands} first, then you.

\begin{careful}
A run's \code{-{}-yolo}, \code{-{}-judge} or \code{-{}-unattended} is a
property of \emph{that run}. A daemon left running never carries it into a
session someone opens later from a browser: that session's confirms follow
\code{[approval]}. A run whose requests nobody answers by hand cannot be
started or resumed from a paired device.
\end{careful}

\subsection{Unattended runs}

\begin{sh}
nodwyn run --task task.md --repo . --max-turns 20 --budget-wall 2h --unattended
\end{sh}

Nobody answers an unattended run, so nothing in it waits for anybody. Tool
confirms are declined unless a standing approval covers them; writes in the
folder are one by default, so an unattended run writes as a chat does
(with \code{-{}-yolo} too, ordinary confirms are answered yes; security
confirms are declined either way). A soft budget checkpoint, if you set one
lower than its hard limit, goes on to the hard limit, the spend limit's
included: give the run \code{-{}-budget-cost 5} and it ends \code{limited}
once its known spend reaches \$5, kept on resume.
The run checkpoint is answered by the plan: after a turn that marked a step
done with \code{plan-step}, the next turn starts; after a turn that marked
none, the run stops and says so, for you to read the turn and
\code{nodwyn run resume -{}-unattended SESSION}. A plain resume is
supervised: how a run goes on is said again at every resume.

\subsection{Judged approval}

With \code{-{}-judge} instead of \code{-{}-yolo} (the two exclude each
other), the judge answers the run's ordinary confirms, in an unattended run
or a supervised one. It needs a judge: a \code{jev-key} stored and
\code{jev.content} other than \code{none} (Chapter~\ref{ch:keys}); without
one the run does not start. The judge is shown the task, the plan and the
call (the tool, why it needs approval, the path it writes or the command
line it runs; at \code{jev.content = "full"} also what it would write or
send) and asked whether to allow it; the call runs only if the probability
of allow is at least \code{jev.approve\_threshold} (0.8) in your config.
The answer is for that call only, never a standing one, and the log notes
each one with the judge's probability. Security confirms are declined
without asking, and a failed judge is a no.

The decision is made in Go, so no Scheme, a repository's included, can
approve a call the judge refused. A project may define
\code{(approve-refuse? call)} in \file{.nodwyn/init.scm}: it is asked about
a call the judge allowed, with \code{call} an association list of
\code{task}, \code{plan}, \code{prim}, \code{reasons}, \code{path},
\code{argv}, \code{args} and \code{detail}, and true, or an error,
declines it. It can only refuse more.

A chat asks the same judge, in the same way, about the commands no
standing approval covers (\code{[approval] commands = "judge"}, the
default; Chapter~\ref{ch:primitives}). There a refusal, or a failed judge,
is not a no: it is a question for you, and the prompt says what the judge
answered. Security confirms, writes, calls that carry a credential and
uploads are never put to it there; they go to you.

\begin{puzzle}[ 17: plan only for big asks]
Write a \code{nodwyn-turn} that records your house plan only when the
request is longer than 200 characters or mentions ``refactor'', and
otherwise goes straight to \code{agent-turn}. What happens in a supervised
run when the request is short?\ans{17}
\end{puzzle}
